\documentclass[b4paper,7pt,twocolumn,landscape]{jsarticle}
\usepackage{info_literacy_exam}

\begin{document}

\twocolumn[
\begin{center}
\large{\textbf{情報リテラシー 2026年度 前期中間試験 問題用紙}}
\end{center}

\noindent
出席番号：\underline{\hspace{2cm}} \quad 氏名：\underline{\hspace{6cm}}

\vspace{0.5mm}
]

\subsubsection*{注意事項}
{
\begin{itemize}
  \item マークシートの学年は、\textbf{1}を記入・マークしてください。
  \item マークシートのクラスは、M:01、E:02、I:03、C:04、A:05の該当する数字を記入・マークしてください。
  \item マークシートの番号は、出席番号を\textbf{3桁}に変換して、右詰めにして、記入・マークしてください。
  \item 問題は全部で\textbf{50問}あります。\textbf{100点満点}（各2点）です。
  \item 解答は\textbf{マークシート}の解答欄に\textbf{正確に濃く記入}してください。
  \item すべて4択問題です。各設問について最も適切な選択肢を1つ選び、マークしてください。
  \item 試験時間は\textbf{50分}です。\textbf{持ち込み不可}（教科書・ノート・電子機器など使用できません）。
  \item 到達目標1：問1$\sim$問50
\end{itemize}
}

% ============================================================
% 【第1回】学内ネットワークと学習環境の基礎（問1〜問6）
% ============================================================
\subsubsection*{【第1回】学内ネットワークと学習環境の基礎（問1〜問6）}

\begin{enumerate}
\item 学内LANを通じてアクセスできるシステムとして、授業で挙げられて\textbf{いない}ものはどれか。\\
\choices{WebClass|Gmail|Teams|LINE}

\item 共用PCで作業を終了した後に必ず行うべきこととして最も適切なものはどれか。\\
\choices{電源を切る|ログアウトする|ファイルを削除する|ブラウザを閉じる}

\item Gmailのセキュリティを向上させるために設定すべき認証方式はどれか。\\
\choices{パスワード認証|指紋認証|多要素認証（2段階認証）|QRコード認証}

\item 本校のパスワードポリシーにおけるパスワードの最小文字数として正しいものはどれか。\\
\choices{8文字|12文字|16文字|20文字}

\item パスワードを変更するためにアクセスするシステムはどれか。\\
\choices{WebClass|Teams|AXIOLE（アクシオーレ）|Gmail}

\item 学内システムへのログインに関する説明として正しいものはどれか。\\
\choices{各システムごとに異なるアカウントが必要|学校のアカウントでログインする|個人のGoogleアカウントを使用する|SNSアカウントで連携ログインする}
\end{enumerate}

% ============================================================
% 【第2回】TeamsとFormsの基本操作（問7〜問13）
% ============================================================
\subsubsection*{【第2回】TeamsとFormsの基本操作（問7〜問13）}

\begin{enumerate}
\setcounter{enumi}{6}
\item 個人Microsoftアカウント作成時に取得するメールアドレスのドメインはどれか。\\
\choices{@gmail.com|@outlook.jp|@yahoo.co.jp|@microsoft.com}

\item Teamsの主な役割として最も適切なものはどれか。\\
\choices{回答の収集・集計|コミュニケーション|アンケート作成|ファイル保存専用}

\item Formsの主な用途として最も適切なものはどれか。\\
\choices{チャット・通話|ファイル共有|アンケート・テスト作成|ビデオ会議}

\item Teamsのチャットとチャネル投稿の違いとして正しいものはどれか。\\
\choices{チャットは全体への発信である|チャットは1対1やグループの会話、チャネル投稿はクラス全体への発信|両方とも同じ機能である|チャネル投稿はファイル共有専用である}

\item クラス全員から授業の感想を集めたい場合、最も適切なツールはどれか。\\
\choices{Teams|Forms|OneDrive|Gmail}

\item デジタルコミュニケーションのマナーとして適切なものはどれか。\\
\choices{@全員を頻繁に使って連絡する|投稿は簡潔・丁寧にし、スレッド返信を活用する|リアクションだけで返信し言葉は添えない|件名を「質問」だけにしてメールを送る}

\item Teams会議に参加する際、「今すぐ参加」をクリックする前に確認すべきことはどれか。\\
\choices{チャット欄にメッセージを送信する|カメラ・マイクのオフを確認する|退出ボタンの位置を確認する|画面共有の設定をする}
\end{enumerate}

% ============================================================
% 【第3回】情報とメディアの特性（問14〜問21）
% ============================================================
\subsubsection*{【第3回】情報とメディアの特性（問14〜問21）}

\begin{enumerate}
\setcounter{enumi}{13}
\item 「データ」の説明として最も適切なものはどれか。\\
\choices{受け手の行動や判断に影響を与えるもの|事実・観測結果をそのまま記録したもの|意味や解釈が加わったもの|分析結果をまとめたもの}

\item 「25」というデータに文脈を加えて「情報」にした例として正しいものはどれか。\\
\choices{25|二十五|気温25℃|0x19}

\item デジタル情報の「複製性」の説明として正しいものはどれか。\\
\choices{情報は瞬時に広範囲に広がる|コストをかけずに同じ情報を大量に複製できる|デジタル情報は削除しても完全には消えない|情報は形式によって伝わりやすさが変わる}

\item 一度発信した情報は回収が難しいという特性を何というか。\\
\choices{複製性|伝播性|残存性|形式性}

\item 「デジタルタトゥー」と最も関連が深い情報の特性はどれか。\\
\choices{複製性|伝播性|残存性|形式性}

\item 情報の評価5視点に含まれ\textbf{ない}ものはどれか。\\
\choices{正確性|網羅性|経済性|客観性}

\item 情報を保存するための媒体を何というか。\\
\choices{表現メディア|伝達メディア|記録メディア|通信メディア}

\item 伝達メディアとしてのインターネットの特徴の組合せとして正しいものはどれか。\\
\choices{高速・世界規模・双方向性あり|リアルタイム・放送エリア・双方向性なし|低速・配布エリア・双方向性なし|リアルタイム・通信可能エリア・双方向性あり}
\end{enumerate}

% ============================================================
% 【第4回】課題の見つけ方と解決プロセス（問22〜問29）
% ============================================================
\subsubsection*{【第4回】課題の見つけ方と解決プロセス（問22〜問29）}

\begin{enumerate}
\setcounter{enumi}{21}
\item 問題解決の5ステップの最初のステップはどれか。\\
\choices{情報の収集|問題の発見|解決策の設計|情報の分析}

\newpage

\item PDCAサイクルの「C」は何を意味するか。\\
\choices{Create（作成）|Check（評価）|Change（変更）|Control（管理）}

\item PDCAサイクルの説明として正しいものはどれか。\\
\choices{1回実行すれば終了する|繰り返すことで継続的に改善する|計画と実行の2段階で構成される|評価は不要である}

\item AND検索の説明として正しいものはどれか。\\
\choices{いずれかのキーワードを含むページを検索|複数のキーワードをすべて含むページを検索|特定のキーワードを除いて検索|完全一致するページのみ検索}

\item 自分でアンケートや観察によって収集した情報を何というか。\\
\choices{二次情報|一次情報|統計情報|公開情報}

\item ブレーンストーミングのルールとして正しいものはどれか。\\
\choices{質を重視し厳選する|他人の意見を批判して改善する|批判しない・自由に発言・質より量|1人で静かに考える}

\item KJ法の手順として正しい順序はどれか。\\
\choices{グループ化→カード作成→文章化→図解化|カード作成→グループ化→図解化→文章化|図解化→カード作成→グループ化→文章化|文章化→図解化→グループ化→カード作成}

\item 中心にテーマを置き、関連する語句を放射状に広げて整理する手法はどれか。\\
\choices{ブレーンストーミング|KJ法|マインドマップ|PDCAサイクル}
\end{enumerate}

% ============================================================
% 【第5回】知的財産と個人情報の基礎（問30〜問37）
% ============================================================
\subsubsection*{【第5回】知的財産と個人情報の基礎（問30〜問37）}

\begin{enumerate}
\setcounter{enumi}{29}
\item 著作権の発生に関する説明として正しいものはどれか。\\
\choices{特許庁に出願して初めて発生する|著作物を販売した時点で発生する|著作物が創作された時点で自動的に発生する|著作権協会に登録して初めて発生する}

\item 日本における著作権の保護期間として正しいものはどれか。\\
\choices{著作者の死後50年|著作者の死後70年|著作物の公表から50年|著作物の公表から70年}

\item 製品のデザイン（外観）を保護する知的財産権はどれか。\\
\choices{特許権|実用新案権|意匠権|商標権}

\item 著作権法上、許可される行為として正しいものはどれか。\\
\choices{他人の文章を出典なしに自分のレポートに使う|音楽を無断でコピーして友人に配布する|私的使用のために著作物を複製する|他人の画像を無断でWebに転載する}

\item 正しい引用の条件に含まれ\textbf{ない}ものはどれか。\\
\choices{引用する必然性がある|引用部分が自分の文章より多い|引用部分と自分の文章が明確に区別されている|出典を明記する}

\item 個人情報の説明として正しいものはどれか。\\
\choices{企業の財務情報のことである|生存する個人を特定できる情報のことである|インターネット上の情報すべてを指す|匿名のアンケート結果のことである}

\item プライバシー権の説明として最も適切なものはどれか。\\
\choices{自分の発明を保護する権利|自分の顔や姿を無断で撮影されない権利|自分に関する情報を他人に知られず自己決定できる権利|自分の著作物を無断で使用されない権利}

\item SNSに写真を投稿する際のリスクとして適切で\textbf{ない}ものはどれか。\\
\choices{写真の背景から住所が特定される|位置情報（Exif）から撮影場所が判明する|写真のファイルサイズから個人が特定される|他人が映り込んだ写真の公開が肖像権侵害になる}
\end{enumerate}

% ============================================================
% 【第6回】(問38〜問43) — PCの設定1：基本設定とネットワーク設定
% 正解: 38→3, 39→3, 40→2, 41→1, 42→4, 43→2
% ============================================================
\subsubsection*{【第6回】PCの設定1：基本設定とネットワーク設定（問38〜問43）}

\begin{enumerate}
\setcounter{enumi}{37}
\item PCを動作させるための基本ソフトウェアを何と呼ぶか。\\
\choices{アプリケーション|ブラウザ|OS（オペレーティングシステム）|ドライバ}

\item Windowsの初期セットアップでサインインに使用するアカウントとして，本授業で指定されているものはどれか。\\
\choices{Googleアカウント|学校のメールアカウント|個人Microsoftアカウント|ローカルアカウント}

\item WindowsのPINに関する説明として正しいものはどれか。\\
\choices{英字のみで構成する必要がある|4桁以上の数字で設定する|メールアドレスと同じものを設定する|一度設定すると変更できない}

\item Windowsのデバイス名を設定する際のルールとして正しいものはどれか。\\
\choices{英数字とハイフンで設定する|日本語の名前を使用する|自動的に割り当てられ変更できない|スペースを含めて自由に設定できる}

\item Windows Updateの主な目的として最も適切なものはどれか。\\
\choices{デスクトップの壁紙を変更する|不要なアプリケーションを削除する|インターネットの接続速度を向上させる|セキュリティの修正や機能の改善を適用する}

\item 学内Wi-Fiに接続する際に必要な情報の組み合わせとして正しいものはどれか。\\
\choices{MicrosoftアカウントとPIN|SSIDとユーザID・パスワード|LANケーブルとACアダプタ|メールアドレスとデバイス名}
\end{enumerate}

% ============================================================
% 【第7回】(問44〜問50) — OfficeとTeamsのインストール
% 正解: 44→4, 45→3, 46→1, 47→2, 48→3, 49→4, 50→3
% ============================================================
\subsubsection*{【第7回】OfficeとTeamsのインストール（問44〜問50）}

\begin{enumerate}
\setcounter{enumi}{43}
\item Microsoft 365に含まれるアプリケーションの組み合わせとして正しいものはどれか。\\
\choices{Word・Safari・Teams|Chrome・Excel・Outlook|Word・Excel・Photoshop|Word・Excel・PowerPoint}

\item Microsoft 365をインストールする際にアクセスするWebサイトはどれか。\\
\choices{www.microsoft.com|teams.microsoft.com|portal.office.com|outlook.office.com}

\item 本授業でMicrosoft 365をインストールする際に使用するアカウントはどれか。\\
\choices{学校のMicrosoftアカウント|個人Microsoftアカウント|Googleアカウント|ローカルアカウント}

\item OneDriveの説明として正しいものはどれか。\\
\choices{Webブラウザの一種である|クラウド上にファイルを保存するサービスである|ウイルス対策ソフトウェアである|メールを送受信するアプリケーションである}

\item Microsoft Teamsの主な用途として最も適切なものはどれか。\\
\choices{画像の編集・加工|ウイルスの検出・駆除|チームでのコミュニケーションや共同作業|ハードウェアの管理}

\item Microsoft 365をインストールする前にWi-Fiをオフにして有線LANに接続する理由として最も適切なものはどれか。\\
\choices{Wi-Fiではインストールができないため|Wi-Fiの電波がPCに悪影響を与えるため|有線LANの方がセキュリティが低いため|有線LANの方が安定した通信でインストールできるため}

\item Microsoft 365のインストール中にトラブルが発生した場合の対処として，最も不適切なものはどれか。\\
\choices{ブラウザを変えて再試行する|有線LAN接続を確認する|自己判断でアンインストールする|担当教員・技術職員に相談する}
\end{enumerate}

\end{document}
